%----------------------------------------------------------------
%
%  File    :  survey-relatedwork.tex
%
%  Author  :   Lukas Auer, David Heidinger, Nina Tschikof and Christina Vogel, TU Graz, Austria
% 
%  Created :  06 May 2026
% 
%  Changed :  06 May 2026
% 
%----------------------------------------------------------------


\chapter{Concluding Remarks}
\label{chap:Conclusion}

Data visualisations are powerful tools for communicating complex information, but they can also be used in misleading ways. This paper examined a range of dishonest chart techniques that distort the interpretation of data through manipulations of selection, aggregation, scaling, visual encoding, and uncertainty representation. Although these techniques often preserve the underlying data values, they can significantly alter how viewers perceive patterns, trends, and relationships.

To better understand these forms of distortion, this paper proposed a taxonomy consisting of nine common dishonest chart techniques. To support practical understanding, interactive demonstrations were developed that allow users to directly experience how visual manipulations influence perception. 

Overall, this work aims to contribute to a more critical understanding of data visualisations and to promote more transparent and responsible visual communication. As visualisations continue to shape public opinion and decision-making, the ability to recognise dishonest chart techniques becomes increasingly important for both creators and consumers of visual information.

